\section{Independent audit of the 244-state fold}
\label{sec:independent-local-audit}

This section gives an independent proof review of candidate~1.  Its checker,
\texttt{verify\_local\_stuttering.py}, reads only the two literal transition
tables.  It does not import the searcher, the decoder compiler, a register
interpreter, or the symbolic frame checker.  The input Routing245 table is
pinned by SHA-256
\begin{center}\small\ttfamily
f8dbeacdad5fd89c57e3f240781324b09110be8e0ccdb838a9c710fb28d2a267
\end{center}
and the verified candidate table has SHA-256
\begin{center}\small\ttfamily
d603bc3793ac9373f510a05d4e898f4ee7b4cac492821406a665554dca5c6daf.
\end{center}
These are the old-label tables: their common initial state is $95$, and
their common target-only halt state is $254$.

\subsection{What is checked independently}

The old machine $M$ has $245$ working states.  The new machine $N$ deletes
state $122$ and changes precisely four retained rows:
\[
\begin{array}{c|c|c}
\text{row}&M&N\\ \hline
(92,1)&(1,L,24)&(1,R,66)\\
(104,1)&(1,R,122)&(1,R,92)\\
(106,1)&(1,R,122)&(1,R,92)\\
(167,1)&(1,L,92)&(1,L,151).
\end{array}
\]
Both rows at $122$ are deleted.  The literal inventories therefore give
$244$ working states and $488$ rows for $N$.

At comparison boundaries, use identical configurations $(q,h,T)$ of the
two machines, with $q\notin\{92,122\}$.  Impose the local condition
\[
\begin{aligned}
q=68&\ \Longrightarrow\ T(h-1)=1,\\
q=81&\ \Longrightarrow\
  \bigl(T(h-1)=0\ \lor\ [T(h-2)=T(h-1)=1]\bigr),\\
q=167&\ \Longrightarrow\ T(h-1)=1.
\end{aligned}
\tag{I}
\]
The condition is vacuous at state $95$, for every tape and head position.
The complete incoming-edge lists for its three constrained states are the
same in both tables:
\[
\begin{array}{c|l}
68&(52,1)\mapsto(1,R,68),\quad(123,1)\mapsto(1,R,68)\\
81&(68,1)\mapsto(1,R,81),\quad(123,0)\mapsto(0,R,81)\\
167&(145,1)\mapsto(1,R,167),\quad(148,1)\mapsto(1,R,167).
\end{array}
\]
The first and third rows establish the required left-hand $1$ by a literal
write.  In the middle row, the entry from $123$ writes the left-hand $0$;
the entry from $68$ writes the nearest $1$ and inherits the next $1$ from
condition~(I).  This proves preservation of~(I) by ordinary comparison
steps.  No condition concerning a register value or a source instruction
is used.

There are $482$ ordinary rows: every row out of a boundary state except
$(81,1),(104,1),(106,1),(167,1)$.  The checker verifies that all $482$ rows
are identical and that none enters $92$ or $122$.  All other boundary
exits are covered by the following seven local calculations.  Coordinates
are relative to the starting head; unspecified cells are arbitrary.
\[
\begin{array}{c|l|c|c|c|l}
q&\text{fixed input cells}&\text{old steps}&\text{new steps}
  &\text{endpoint }(q',h')&\text{changed cells}\\ \hline
104&T(0)=1,\ T(1)=0&2&4&(25,2)&\text{none}\\
104&T(0)=1,\ T(1)=1&2&2&(66,2)&\text{none}\\
106&T(0)=1,\ T(1)=0&2&4&(25,2)&\text{none}\\
106&T(0)=1,\ T(1)=1&2&2&(66,2)&\text{none}\\
167&T(-1)=T(0)=1&2&2&(24,-2)&\text{none}\\
81&T(-1)=0,\ T(0)=1&2&2&(107,-2)&\text{none}\\
81&T(-2)=T(-1)=T(0)=1&6&4&(2,-2)&T(-1),T(0):1\mapsto0.
\end{array}
\tag{L}
\]
In each row, the two executions finish in the same state, at the same
head position, with the same entire tape.  The checker rejects a trace
that reads any cell outside its fixed input window.  Thus this is an
exact verification for arbitrary exterior tape contexts, rather than
testing a finite collection of complete tapes.

The two input cases for each of $104$ and $106$ exhaust the binary value
of $T(1)$.  The two cases at $81$ exhaust condition~(I).  The case at $167$
uses precisely its condition in~(I).  Every endpoint in~(L) satisfies~(I)
vacuously and is a comparison boundary.  This establishes a positive
stuttering relation without needing comparisons at intermediate states
$107$, $125$, or $151$ visited by one of the short detours.

\subsection{Halting equivalence and time accounting}

\paragraph{Theorem.}
For every function $T:\mathbb Z\to\{0,1\}$ and every head position $h$,
$M$ started in $(95,h,T)$ halts if and only if $N$ started in the same
configuration halts.

\paragraph{Proof.}
Start with the common state $95$, where~(I) holds.  At every boundary,
take an ordinary common step or the applicable row of~(L).  The endpoints
coincide exactly, remain boundaries, and satisfy~(I).  Induction constructs
paired boundary configurations.  Every block takes at least one step in
each machine; its lengths are at most six old steps and four new steps.
No exceptional block contains the halt state.  Ordinary halting steps
are identical.  Therefore a finite halt in one run occurs at the common
endpoint of a finite sequence of blocks in both runs.  If the block
sequence is infinite, its two positive clocks are unbounded, so both
runs are nonhalting.\hfill$\square$

The argument is independent of the old machine's blank initialization.
Consequently it transfers the old machine's consistency calibration to
the exact new table by composition.  A kernel-checked endpoint for the
new literal table is not claimed by this standalone Python checker.

\paragraph{Corollary (exact clock difference).}
If the common computation halts, let $A$ count its comparison blocks
that start at $104$ or $106$ with the next bit $0$, and let $B$ count its
blocks that start at $81$ with the two preceding bits $11$.  If
$t_M,t_N$ are the two complete halting times, then
\[
 t_N=t_M+2A-2B,
 \qquad \frac23 t_M\leq t_N\leq 2t_M.
\]
Indeed, exactly the $A$ blocks add two steps and exactly the $B$ blocks
remove two steps.  Every individual block satisfies
$2\ell_M\leq3\ell_N$ and $\ell_N\leq2\ell_M$; summing gives the bounds.
These are statements about complete halting runs, not about alignment
of arbitrary individual microsteps.

For the blank tape, both halting predicates are arithmetical.  The
proof uses finite case checks and induction over coded finite
computations, so its blank-tape specialization can be formalized in PA.
There is no reliance on second-order quantification over all tapes for
that arithmetical specialization.  To see explicitly that the reverse
direction also uses finite reasoning, an alleged new halt within $t$
steps can be compared against an old prefix of $6t$ steps.  Unless the
old machine has already halted, this prefix supplies at least $t$
positive comparison blocks, whose new counterpart reaches time at least
$t$ without a halt, a contradiction.

\subsection{Audit conclusion}

The independent checker passed all seven exceptional local identities
and all $482$ ordinary rows.  Its local executions contain $18$ old and
$20$ new literal transitions in total.  Its complete output is
\texttt{local-stuttering-audit.json}.  The checks support a direct theorem
between the two exact Turing machines; they do not depend on the larger
register-frame certificate or on bounded numerical experiments.
